\documentclass{article}
\usepackage{geometry}
\geometry{a4paper, margin=1in}
\usepackage{graphicx}
\usepackage{hyperref}

\title{Practical Machine Learning and Deep Learning\\D1.2 Progress Submission\\
Team "Magic Photo"}
\author{Maksim Ilin \href{mailto:m.ilin@innopolis.university}{m.ilin@innopolis.university},\\ Rail Sabirov \href{mailto:rai.sabirov@innopolis.university}{rai.sabirov@innopolis.university},\\ Ivan Ilyichev \href{mailto:i.ilyichev@innopolis.university}{i.ilyichev@innopolis.university}}

\date{\today}

\begin{document}

\maketitle

\section{Project Topic}
\par This project objective is to enhance existing solutions for image restoration from different corruptions and coloring of grayscale images using machine learning and deep learning techniques. The project will focus on developing a user-friendly application that allows users to upload images, select the processing option they want to apply, and receive the processed image as output. Additionally, we plan to collect users' feedback for future product improvements and keep track of software lifecycle. The application will leverage pre-trained models and fine-tune them. 

\section{Link to the Repository}
\par The link: \url{https://github.com/ivaniinno/Image-Restoration-and-Coloring/tree/dev}. For now we store code in \textbf{dev} branch with \textbf{data\_augmentation} folder, \textbf{docs} folder with project documentation, \textbf{models\_testing} folder for trying baseline model implementations, and in \textbf{models\_finetuning} we have a template for finetuning. Branch \textbf{issue-2} contains exploratory data analysis.


\section{Dataset and Exploratory Data Analysis}

\par Out datasets is composed from three sources:

\begin{itemize}
    \item Nature Dataset: \url{https://huggingface.co/datasets/mertcobanov/nature-dataset} (\textbf{nature} category in our dataset)
    \item European Cities 1M dataset: \url{http://image.ntua.gr/iva/datasets/ec1m/index.html} (\textbf{city} category in our dataset)
    \item CelebA Dataset: \url{https://www.kaggle.com/datasets/jessicali9530/celeba-dataset} (\textbf{faces} category in our dataset)
\end{itemize}

\par One can find image examples in Figure \ref{fig:example_images}

\begin{figure}[!hbt]
    \centering
    \includegraphics[width=1\textwidth]{0.png}
    \caption{Images Examples from Each Category}
    \label{fig:example_images}
\end{figure}

\par To balance category distribution, we randomly sampled 50000 images from each source.
\par Results of data exploration:
\begin{enumerate}
    \item 56 duplicates were found and removed
    \item 2743 gray images were deleted because they are not suitable for the colorization task (we cannot compare the result with the original image) as shown on Figure \ref{fig:gray_images}.
    \item \textbf{faces} category has resolution of 178x218, \textbf{nature} category has resolution of 512x512, and \textbf{city} category resolutions varies as Figures \ref{fig:width_and_height} and \ref{fig:aspect_distribution} demonstrate. Figure \ref{fig:aspect_distribution} shows that most of the images have aspect ratio 1-2. Images with ratios more than 6 were removed as extreme cases (in total 24 of them).
    \item Images were divided into classes via ResNet (for feature extraction), and PCA was used for visualization (see Figure \ref{fig:images_visualization}). In genereal, classes are well separable, however, images from the \textbf{city} category slightly overlap with two other clusters, which indicates the presence of common visual patterns between the \textbf{city} category and the other two classes.
    \item Average brightness distribution is as shown on Figure \ref{fig:brightness_distibution}. 1457 outliers ($\approx$ 1\%) were found (underexposed and overexposed images) (see Figure \ref{fig:underexposed_overexposed_images}). These images were saved and will be used for training the model on hard cases.
    \item As a result, 147177 images remain from 150000:
    \begin{itemize}
        \item city - 48112
        \item faces - 49985
        \item nature - 49080
    \end{itemize}
    The final classes proportion seems to be appropriate.

\end{enumerate}

\begin{figure}[!hbt]
    \centering
    \includegraphics[width=1\textwidth]{1.png}
    \caption{Amount of Gray Images in Each Category}
    \label{fig:gray_images}
\end{figure}

\begin{figure}[!hbt]
    \centering
    \includegraphics[width=1\textwidth]{2.png}
    \caption{Images Height and Width Distribution}
    \label{fig:width_and_height}
\end{figure}

\begin{figure}[!hbt]
    \centering
    \includegraphics[width=1\textwidth]{3.png}
    \caption{Images Height and Width Comparison}
    \label{fig:aspect_distribution}
\end{figure}

\begin{figure}[!hbt]
    \centering
    \includegraphics[width=1\textwidth]{4.png}
    \caption{Image Features and Classes Visualization}
    \label{fig:images_visualization}
\end{figure}

\begin{figure}[!hbt]
    \centering
    \includegraphics[width=1\textwidth]{5.png}
    \caption{Images Brightness Distribution}
    \label{fig:brightness_distibution}
\end{figure}

\begin{figure}[!hbt]
    \centering
    \includegraphics[width=1\textwidth]{6.png}
    \caption{Underexposed and Overexposed Images Counts}
    \label{fig:underexposed_overexposed_images}
\end{figure}

\section{Data Augmentation}
\par Several data augmentation techniques were implemented to increase dataset size and variability.
\par Degradation Types include:

\begin{enumerate}
\item Noise: Gaussian, salt-pepper, speckle
\item Blur: Gaussian, median, motion blur
\item Resolution: Downscaling with upscaling artifacts
\item Compression: JPEG artifacts (quality 20-30)
\item Physical damage: Scratches and dust spots
\item Color: Grayscale conversion
\end{enumerate}

\par Additionaly, mentioned degradations were combined to create more complex scenarios as follows:
\begin{enumerate}
\item \textbf{combined\_restoration}: Dust\_and\_Scratch/None + noise(Gaussian/SP/Speckle) + low resolution/compression
\item \textbf{combined\_superres}: Blur(Gaussian/Median/Motion\ blur) + resolution reduction/compression
\item \textbf{combined\_colorization}: Grayscale + Noise(Gaussian/SP/Speckle)/Blur(Gaussian/Median/Motion\ blur)/None
\end{enumerate}

\section{Model Testing}
\par Real-ESRGAN and GPFGAN models were tested out-of-the-box on selected images for future comparisons after fine-tuning. The notebook is located at \textbf{models\_testing/restoration\_superres/restorationtest}. Preferable link to open the notebook: \url{https://www.kaggle.com/code/railsabirov/restorationtest}.

\par The results are on Figure \ref{fig:model_testing_results}.

\begin{figure}[!hbt]
    \centering
    \includegraphics[width=1\textwidth]{7.png}
    \caption{Real-ESRGAN and GPFGAN Models Results on Different Degradations}
    \label{fig:model_testing_results}
\end{figure}

\par As one can see, on average PSNR shows 24-27 dB, and SSIM shows 0.74-0.81. Such metrics show that the restored image is quite good in quality. The combination of Real-ESRGAN + GPFGAN works best for restoring images with faces, then comes the urban landscape, and it works worst with the nature category, since there is a lot of noise initially in the image, which the model smooths out.

\section{Future Work}
\par In the closest perspective we are planning from highest to lwoest priority (first 2 points are crucial and to be done first):

\begin{itemize}
    \item Fine-tune Real-ESRGAN and GPFGAN models.
    \item Investigate other models for possible consideration of them.
    \item Start developing simple backend/frontend for interaction with model.
\end{itemize}

\end{document}